\section{Mechanization}\label{sec:mech}

The development has nine CryptoVerif models: two protocol models; five KEM
models, namely two variants each of \file{kem_xwing.m4.ocv} and
\file{kem_bbfgs.m4.ocv} and the standalone \file{kem_bbfgs_x25519_rom.ocv};
and two variants of the composite-signature model. Together they come
to about 1,400 lines, including comments, and \texttt{make} checks all of them
in under 20~seconds on two cores. \Cref{tab:models} lists them.

\begin{table}[t]
\centering
\small
\begin{tabularx}{\linewidth}{@{}lLll@{}}
\toprule
Model & Proves & Library & Queries\\
\midrule
\file{hybrid_ake.ocv} & \cref{thm:auth,thm:secrecy}; key uniqueness & pq & 5\\
\file{hybrid_ake_ephleak.ocv} & the same with static ephemeral reveal & pq & 5\\
\file{composite_sig.m4.ocv} (2 variants) & \cref{lem:sig}, each component & pq & 1 + 1\\
\file{kem_xwing.m4.ocv}, ML-KEM variant & \cref{lem:kem}, X-Wing, ML-KEM leg & pq & 1\\
\file{kem_xwing.m4.ocv}, X25519 variant & \cref{lem:kem}, X-Wing, X25519 leg & default & 1\\
\file{kem_bbfgs.m4.ocv}, ML-KEM variant & \cref{lem:kem}, ML-KEM-1024, ML-KEM leg & pq & 1\\
\file{kem_bbfgs.m4.ocv}, X25519 variant & \cref{lem:kem}, ML-KEM-1024, X25519 leg (PRF-ODH) & default & 1\\
\file{kem_bbfgs_x25519_rom.ocv} & \cref{lem:kem}, ML-KEM-1024, X25519 leg (ROM) & default & 1\\
\bottomrule
\end{tabularx}
\caption{The nine models. ``pq'' is CryptoVerif's post-quantum-sound library.
The X25519 directions use the default library and are classical by nature.
The composite's EdDSA direction runs with the post-quantum library, but its
assumption, unforgeability of EdDSA, holds only classically.}
\label{tab:models}
\end{table}

\subsection{The protocol model}\label{sec:mech:protocol}

The protocol model abstracts each hybrid primitive as a single primitive: one
KEM satisfying the library's \ident{IND_CCA2_KEM} macro, and one signature
satisfying \ident{UF_CMA_proba_signature}. The transcript hash is a
collision-resistant hash, and the KDF a PRF keyed by the shared secret. The
length-prefixed encodings are injective constructors. The two signed hash
inputs, $\LP$ of seven fields and $\LP(\THH, \sigma_R)$, are declared
disjoint, since a byte string parses as at most one of them. Freshness and
RevealKey are written directly in the initiator's and responders' oracles,
as \texttt{find} conditions over the other sessions' transcripts;
\cref{app:excerpts} shows the initiator. This is the whole proof script:

\begin{lstlisting}
crypto uf_cma_corrupt(sign) rK;      (* composite signature, all honest keys *)
success; simplify;                   (* this proves both agreement queries *)
replace at_nth 6 1 "kdf({_}ssI" "ssA[u_9, u_10]";    (* KEM correctness, *)
replace at_nth 5 1 "kdf({_}ssI" "ssR[u_11, u_12]";   (* in the three branches *)
replace at_nth 2 1 "kdf({_}ssI" "ssA[u_17, u_18]";   (* where a key is tested *)
crypto ind_cca2(encap) seedE;        (* hybrid KEM, all ephemeral keys *)
crypto prf(kdf) *;                   (* KDF keyed by the shared secret *)
success
\end{lstlisting}

After the signature step, every branch in which an initiator's key is tested
knows that its ciphertext came from an honest responder that encapsulated to
this initiator's ephemeral key. The three \texttt{replace} steps use KEM
correctness to replace the initiator's decapsulated secret with the
responder's encapsulated one. Without them, the IND-CCA2 step cannot
recognise the decapsulation as one of a challenge ciphertext. CryptoVerif
proves each equality, so a \texttt{replace} that targets the wrong
occurrence fails rather than proving something false. The cost is
brittleness: the script names variables that CryptoVerif generates
(\ident{u_9}, \ldots) and counts occurrences, so it is tied to
CryptoVerif~2.13 and to the model text.

\subsection{The KEM models}\label{sec:mech:kem}

Each KEM model proves an equivalence between the real hybrid KEM and the
ideal KEM of CryptoVerif's \ident{IND_CCA2_KEM} macro. In the ideal
process, encapsulation to an honest public key returns a fresh random key
and records the pair in a table, and decapsulation of a recorded ciphertext
returns the recorded key. A setup oracle creates each user from a pair of
seeds, one for ML-KEM and one for X25519, and an \texttt{Oleak} oracle hands
the adversary the broken leg's secret: this is what ``the other leg is fully
broken'' means. Every proof starts by inserting into the real process the
two case distinctions that the ideal one makes: is the encapsulation for an
honest key, and is the decapsulated ciphertext one of that user's challenge
ciphertexts. It then uses ML-KEM's correctness on the second, through the
library's perfect-correctness equation. After that, the ML-KEM direction
applies IND-CCA2 of ML-KEM and then the PRF assumption on the combiner, keyed
by the now uniformly random $\ssv_M$. The X25519 direction applies the random
oracle and then strong Diffie--Hellman, except in the PRF-ODH model, which
applies the PRF-ODH assumption directly. In all cases
\ident{start_from_other_end} transforms the ideal process to the same
final game.

\paragraph{C2PRI.} In X-Wing's X25519 direction, ML-KEM has no security
assumption except C2PRI (Definition~7 of~\cite{XWing}): for an honestly
generated key pair and an honest encapsulation $(k_0, c_0)$, an adversary
given the decapsulation key, $k_0$ and $c_0$ finds no $c \ne c_0$ that
decapsulates to $k_0$. We state it as a CryptoVerif collision statement:
\begin{lstlisting}
collision s <-R mseed; r <-R mencseed; forall c: mct;
   return(mdecap(c, mskgen(s)) = minjbot(1-proj-mpair(mencap_r(mpkgen(s), r))))
   <=(Pbind(time))=> return(c = 2-proj-mpair(mencap_r(mpkgen(s), r))).
\end{lstlisting}
($\cvp{Pbind}$ is $\cvp{Pc2pri}$ in \cref{lem:kem}.) For an honestly
generated key seed $s$ and honest encapsulation coins $r$, the statement
replaces the test $\mathsf{mdecap}(c, \mathit{sk}) = k_0$ by the test
$c = c_0$, at cost $\cvp{Pbind}$. It is stronger than Definition~7 in three
ways. First, the adversary may depend on the key seed, not only on the
decapsulation key, and in the X25519 model it is handed the seed. FIPS~203's
decapsulation key does not determine the seed, so this is a genuinely
stronger, seed-known form of C2PRI, and we assume it. It is the natural form
for X-Wing, whose decapsulation key is the seed. The X-Wing paper's
random-oracle proof that ML-KEM-768 is C2PRI~\cite[Thm.~3]{XWing} argues
only about collisions of the hash functions $G$ and $J$ and does not appear
to use the seed's secrecy, but we have not checked in detail that it
carries over. Second, the adversary may depend on the encapsulation coins
$r$. Third, the statement implies that $c_0$ itself decapsulates to $k_0$.
For ML-KEM these last two add at most $2\delta$, where $\delta$ is the
decapsulation failure probability. Given the decapsulation key and $c_0$,
decryption recovers the encapsulation message, which determines the coins,
unless decryption fails; and $c_0$ decapsulates to $k_0$ unless decryption
fails. The statement concerns honestly generated keys and challenges only. It is neither the malicious-key binding that
ML-KEM lacks~\cite{Sch24} nor needed in that form. The proof applies it three
times, all after the random-oracle step, which gives the factor $3NQ_e$ in
\cref{lem:kem}.

The assumption is needed because of X-Wing's design. X-Wing does not hash
$\ct_M$. If an adversary that knows ML-KEM's secret key could
find a second ciphertext $\ct_M' \ne \ct_M^*$ with the same $\ssv_M$, then
$\ct_M' \,\|\, \ct_X^*$ would decapsulate to the challenge key, which the
decapsulation oracle would then reveal. The ML-KEM-1024 combiner hashes
$\ct_M$, so a second ciphertext changes the hash input. That is why its
X25519-direction proofs need no security assumption on ML-KEM. They use only
its correctness, its implicit rejection, and the negligible chance that
honestly generated keys or ciphertexts collide. \Cref{sec:mutations} checks
both facts.

\paragraph{Two routes for the ML-KEM-1024 profile.} The first X25519-direction
proof for this profile, in \file{kem_bbfgs.m4.ocv}, uses the library's
PRF-ODH macro with HKDF keyed by the Diffie--Hellman value. Here that value
is the second half of the ikm, after the adversary-controlled $\ssv_M$, so
the assumption is a non-standard variant. \file{kem_bbfgs_x25519_rom.ocv}
proves the same statement from strong Diffie--Hellman with HKDF as a random
oracle, which are standard assumptions. The ML-KEM direction of this profile
assumes HKDF is a PRF keyed by the $\ssv_M$ half of $\ssv_M \,\|\, \ssv_X$.
That is one half of the split-key PRF property that the combiners of
Giacon, Heuer and Poettering~\cite{GHP18} rely on, and that Millerjord,
Stebila and Steckel~\cite{MSS26} and Alagic, Bajaj and Kocoglu~\cite{ABK25}
study further: a split-key PRF is a PRF when keyed by either part of its
key, so it implies our assumption, but not conversely. We know of no proof
of even this half for HKDF in this configuration. It is also outside
what~\cite{BBGS23} prove, because their key is the whole message input. CryptoVerif's PQ-TLS model
makes the same kind of assumption for HKDF-Extract over a concatenated hybrid
secret~\cite{BJ24}.

\paragraph{X25519 as a group.} All models treat X25519 as a prime-order group
(the library's \ident{DH_basic} and \texttt{StDH} macros). Clamping, the
cofactor, the twist, low-order points and non-canonical encodings are not
modelled; \cref{sec:limits} discusses what this leaves out. The X-Wing
paper instead uses a ``nominal group''~\cite{ABHKLR21,XWing}, and its
classical bound has a statistical term $2\Delta_N$, with
$\Delta_N < 2^{-126}$ for X25519~\cite{XWing}. So the X25519-direction models make the X-Wing
paper's assumptions up to this difference in how X25519 is modelled. The
32-byte X-Wing seed is modelled as two independent seeds, treating SHAKE-256
as a PRG.

\subsection{The composite-signature model}\label{sec:mech:sig}

The composite model has one honest key pair, a signing oracle, a
verification oracle that records an \texttt{accepted} event for every
verifying pair $(\mathit{ctx}, M)$, and an oracle that hands out the broken
component's secret key. The query states that every accepted
$(\mathit{ctx}, M)$ was signed. The message representative is an injective
function \texttt{mp}$(\mathit{ctx}, \PH(M))$; $\mathsf{Prefix}$ and
$\mathsf{Label}$ are constants of the profile, and $\mathrm{len}(\mathit{ctx})$
is one byte. Unlike the protocol and KEM proofs, this one runs without a
proof script: CryptoVerif finds the two steps, unforgeability of the
unbroken component and collision resistance of $\PH$, by itself.
\Cref{app:excerpts} shows the process.
